\subsection*{Divine help directs active obedience\enspace\properrefs{Coll., Ep., Gosp., Comm., Postcomm.}}

The Collect asks God to guard the Church with perpetual mercy because human
mortality slips without him. Its two movements belong together: divine helps
withdraw the person from harmful things and direct the person toward saving
things. Paul's command to walk in the Spirit gives that direction a moral
form. Christ's command to seek the kingdom first gives it a governing end.
The Communion places the same priority upon those receiving, and the
Postcommunion asks the sacraments to lead them to the effect of perpetual
salvation. Dependence on God sustains a life of action.

Chrysostom makes the character of that action especially clear. In his
commentary on Galatians~5:16--24, bodily substance is not the enemy: the body
serves hearing, preaching, and writing, while Paul's catalogue includes
hostility and faction. In the Matthew homilies, possessing wealth differs
from serving it, and the birds do not abolish labor. Job's stewardship and
Paul's work exemplify the right use of created means (\emph{In Gal.}, ch.~5;
\emph{Hom. in Matt.}~21.1--4; 22.1--3). The flesh--Spirit conflict and the
two-master saying thus meet in the ordering of human life under God; neither
makes material existence evil.

Augustine's \emph{Exposition of Galatians} adds an account of the struggle
within obedience. Desires may trouble a person without receiving consent;
grace makes justice more delightful, while freedom from all such conflict
belongs to resurrection (§§47--51). The Collect's mortal petitioner can
therefore pray honestly without pretending to have become invulnerable.
Charity, patience, peace, and chastity are determinate acts within that
struggle. Aquinas calls fruit delightful action produced by the Spirit and
relates its present sweetness to future beatitude (\emph{Super Gal.}, V,
lects.~6--7). The closing prayer asks that present reception carry the
worshipper toward that final effect.

\subsection*{The Father cares through gifts and helpers\enspace\properrefs{Int., Coll., Grad., Gosp., Off., Postcomm.}}

The Gradual ranks confidence in God above confidence in humans and princes.
The Offertory announces the Lord's angel around those who fear him. Augustine
explains the former by making good people and angels instruments of a goodness
they receive from God; even angelic princes cannot replace its source
(\emph{Enarr. in Ps.}~117.4). Theodoret emphasizes the appointed verses'
contrast between God's power and the mortality, changeability, and brief
authority of human helpers (PG~80:1811--1812). Bellarmine likewise allows
created aid while directing final confidence to God (\emph{Commentary},
p.~370). Help can be real without becoming absolute.

The angelic promise itself has several received applications. Basil treats
the believer's guardian and the moral danger of driving that protection away
(\emph{Hom. in Ps.}~33.5--6). Augustine identifies the angel with Christ, the
messenger of great counsel (Ps.~33, second exposition, §10). Aquinas includes
an angelic minister, Christ, and a prelate (\emph{In Ps.}~33.8); Bellarmine's
explanation begins with a created angel (pp.~91--92). These applications
differ in their immediate referent. Together with the Gradual they teach
reverence for God's protection without making a human office or an angel an
independent object of final trust.

Christ names the Father as the one who feeds birds and clothes grass, knows
his children's needs, and commands the kingdom's priority. Augustine retains
both ordinary provision and severe deprivation within this teaching: Christ's
purse, collections for famine relief, and manual labor are lawful, while
Paul's hunger and nakedness do not disprove providence. God the physician
gives or withholds with perpetual rest in view (\emph{De sermone Domini in monte},
II.17, §§57--58). The Collect's saving direction and the Postcommunion's lasting
end explain why the Mass's confidence exceeds a promise of material safety.

\subsection*{Desire becomes approach, praise, and tasting\enspace\properrefs{Int., Ep., All., Gosp., Off., Comm.}}

The Introit desires God's courts and asks him to regard the face of his
anointed. Augustine interprets that petition as a desire that Christ become
known to all and the better day as eternity (Ps.~83.13--14). Theodoret reads
the face corporately as the saved people who are Christ's body, then relates
the preference for God's house to captives' longing and Christian worship
(PG~80:1543--1544). Bellarmine gives a further, distinct account: the Father's
regard for the Messiah's merits and the excellence of heavenly life
(p.~255). The shared chant supports several received christological
and ecclesial emphases.

The Alleluia's invitation to come and rejoice gives desire a public voice.
Augustine describes approach as restored likeness to God and jubilation as
joy beyond adequate words (Ps.~94.2--3). Hugh of Saint-Cher's exposition places
the opening verse in the Matins Invitatory, where coming involves faith,
obedience, and attentive praise (\emph{Postilla}, Ps.~94:1, ff.~249rb--vb).
Theodoret presents David speaking prophetically in the person of Josiah and
the priests, then extends the hymn to apostles and martyrs
(PG~80:1639--1640). These are interpretations of the psalm's voice and use,
not composition dates.

The Offertory adds experience to invitation: taste and see the Lord's
goodness. Augustine explicitly speaks of Christ's Body and Blood (Ps.~33,
second exposition, §11); Basil joins Christ as true bread with an inward
experience of truth (\emph{Hom. in Ps.}~33.6). Aquinas explains spiritual taste
before sight without explicitly identifying Eucharistic species at this locus
(Ps.~33.9). The sacramental application is therefore particularly direct in
Augustine, while the psalm's invitation remains wider than sensory pleasure.
Christ's command at Communion orders the receiver's desire: God's kingdom
comes first. Chrysostom's citation of the soul's longing in Psalm~83:3 while
explaining Galatians~5:17 already joins the Introit's psalm with the Epistle's
struggle; the desired object matters (\emph{In Gal.}, ch.~5).

\subsection*{The saving offering bears fruit in a people\enspace\properrefs{Coll., Ep., Off., Sec., Comm., Postcomm.}}

The Collect names the Church; the Secret speaks of our sins; the
Postcommunion asks for purification, strength, and a lasting saving effect
for us. The Epistle's catalogue reaches the same communal life through
hatred, quarrels, envy, and divisions, then through charity, peace, patience,
and gentleness. Moral transformation and sacramental grace concern a people
whose members must learn to live together.

The Secret calls the offering a saving victim and asks that it purify sins
and propitiate divine power. Schuster connects this prayer and the
Postcommunion with the strength needed for virtue and the attainment of
eternal salvation (\emph{Sacramentary}, III, pp.~137--138). The Catechism
explains Eucharistic fruit through union with Christ, growth in charity,
preservation from sin, the unity of Christ's body, and commitment to the poor
(CCC~1391--1397). The Galatian virtues give daily conduct to that charity.
The Eucharist's strengthening of love reaches someone whose hunger or
exclusion requires another member's action.

Paul's kingdom warning retains its force beside these prayers. The Spirit's
fruit cannot become a collection of achievements independent of grace;
sacramental reception cannot become indifference to the sins named by the
Apostle. The Catechism places growth in charity and the unity of Christ's body among
Communion's fruits (CCC~1394, 1396). The prayers ask for purification and
strength because the receiving Church needs both.

The orations' transmitted association is independently attested. Wilson's
Gregorian supplement, Gellone~II, Sant Ruf, and Vilabertran place the three
prayers together under Sunday~XV in their respective arrangements. Wilson's
Ottobonian chant annotations, however, distribute the present chants and the
prayer trio differently (pp.~173--174 and p.~173 notes a--c). The prayers have
a history as a set; those witnesses do not establish an unchanged ancient
version of the complete 1962 Mass.

\subsection*{Present provision serves perpetual salvation\enspace\properrefs{Int., Coll., Gosp., Comm., Postcomm.}}

Matthew's discourse specifies food, drink, and clothing before promising that
these things will be added. The Communion's shorter form places the kingdom
first and retains the addition, but its \latin{ómnia} receives its antecedent
from that Gospel context. Reception is not a promise that every ambition will
be supplied. The Collect asks for \latin{salutária}; the Postcommunion asks for
the effect of perpetual salvation. Temporal provision belongs within this
larger direction.

Augustine's account of the divine physician is decisive here. God may give
what serves that direction and withhold what does not, while Paul's need and
his labor remain real (\emph{De sermone Domini}, II.17.57--58). Anthony of
Padua similarly counsels gratitude for gifts and patience under loss in his
Sunday~XV sermon (§§12--17). That sermon treats Matthew~6 within a different
medieval formulary, with Galatians~5:25ff. and a \emph{Miserere} Introit.
Its teaching on use and loss illuminates this Gospel without identifying the
two Sunday arrangements.

Augustine's eternal reading of the Introit's one better day supplies the same
horizon in the language of worship (Ps.~83.14). Earthly days, their work, and
their necessities have value under that end. The closing prayer entrusts both
the receiver's purification and the course ahead to the sacraments' saving
effect. The Mass ends with people who still need food and clothing, still
must work and help their neighbors, and have asked to be led beyond the
limits of all present provision.
